%% chapters/04-metodologi.tex
%% Bab IV — Metodologi
%% Aturan: lihat .cursor/rules/09-bab4-metodologi.mdc
%% Bab ini adalah KONTRIBUSI UTAMA — dijabarkan secara teknis dan reproducible.

\chapter{Metodologi Penelitian}
\label{bab:metodologi}

Bab ini menjabarkan metodologi penelitian secara teknis dan
reproducible: arsitektur usulan \emph{Mamba-xLSTM-Net}, dataset yang
digunakan, ekstraksi fitur, modul interpretabilitas, protokol
pelatihan, rancangan eksperimen, metrik evaluasi, dan lingkungan
komputasi.

\section{Kerangka Penelitian}
\label{subsec:bab4_kerangka}

% Gambar IV.1 — diagram alur penelitian end-to-end
\begin{figure}[ht]
\centering
% \includegraphics[width=0.9\textwidth]{figures/bab4/kerangka.pdf}
\caption{Diagram alur penelitian \emph{end-to-end} dari sinyal getaran
mentah hingga prediksi RUL dan analisis interpretabilitas.}
\label{fig:bab4_kerangka}
\end{figure}

Kerangka penelitian yang diusulkan ditunjukkan pada
Gambar~\ref{fig:bab4_kerangka}. Sinyal getaran mentah dari dataset
(Subbab~\ref{subsec:bab4_dataset}) diproses melalui ekstraksi fitur
multi-domain (Subbab~\ref{subsec:bab4_hi}), kemudian dimasukkan ke
arsitektur Mamba-xLSTM-Net (Subbab~\ref{subsec:bab4_arsitektur}) untuk
menghasilkan prediksi RUL. Modul interpretabilitas
(Subbab~\ref{subsec:bab4_interp}) bekerja secara \emph{post-hoc}
terhadap representasi laten model.

\section{Dataset}
\label{subsec:bab4_dataset}

\subsection{PHM2012 / FEMTO-PRONOSTIA}
\label{subsec:bab4_phm2012}
% Spesifikasi PHM2012 dengan tabel sesuai .cursor/rules/15-domain-rul-bearings.mdc §C.1
\dots

\subsection{XJTU-SY}
\label{subsec:bab4_xjtusy}
% Spesifikasi XJTU-SY dengan tabel sesuai .cursor/rules/15-domain-rul-bearings.mdc §C.2
\dots

\subsection{Praproses dan Segmentasi}
\label{subsec:bab4_praproses}
% Window size, overlap, normalisasi (fit hanya pada train), bearing-wise split
\dots

\section{Ekstraksi Fitur dan \emph{Health Indicator}}
\label{subsec:bab4_hi}

% Tabel 16 fitur (9 time-domain + 7 frequency-domain). Persamaan HI.
\dots

\section{Arsitektur Usulan: Mamba-xLSTM-Net}
\label{subsec:bab4_arsitektur}

% Gambar arsitektur dengan dimensi tensor di tiap blok.
\begin{figure}[ht]
\centering
% \includegraphics[width=0.95\textwidth]{figures/bab4/arsitektur.pdf}
\caption{Arsitektur \emph{Mamba-xLSTM-Net} yang diusulkan. Dimensi
tensor pada tiap blok ditandai dengan notasi $B \times T \times d$.}
\label{fig:bab4_arsitektur}
\end{figure}

\subsection{Modul Encoder Mamba}
\label{subsec:bab4_mamba}
\dots

\subsection{Modul Recurrent xLSTM}
\label{subsec:bab4_xlstm}
\dots

\subsection{Mekanisme Fusi}
\label{subsec:bab4_fusi}
\dots

\subsection{Kepala Regresi RUL}
\label{subsec:bab4_head}
\dots

\section{Modul Interpretabilitas}
\label{subsec:bab4_interp}

\subsection{Top-$k$ Sparse Autoencoder}
\label{subsec:bab4_sae}
\dots

\subsection{Atribusi Fitur: SHAP dan Integrated Gradients}
\label{subsec:bab4_shap_ig}
\dots

\section{Protokol Pelatihan}
\label{subsec:bab4_protokol}

\subsection{Fungsi Kerugian}
\label{subsec:bab4_loss}
\dots

\subsection{Optimizer dan Penjadwal Laju Belajar}
\label{subsec:bab4_optim}
\dots

\subsection{Regularisasi dan \emph{Early Stopping}}
\label{subsec:bab4_regul}
\dots

\section{Rancangan Eksperimen dan Ablasi}
\label{subsec:bab4_eksperimen}

% Daftar terstruktur eksperimen E1, E2, ... lihat aturan
% .cursor/rules/09-bab4-metodologi.mdc §H

\subsection*{Eksperimen E1 — Perbandingan SOTA}
\textbf{Tujuan:} Menjawab RQ1.\\
\textbf{Setup:} \dots\\
\textbf{Output:} Tabel V.1, Gambar V.1.

\subsection*{Eksperimen E2 — Ablasi Komponen Arsitektur}
\textbf{Tujuan:} Memvalidasi kontribusi tiap modul.\\
\textbf{Setup:} \dots

\subsection*{Eksperimen E3 — Ablasi Hyperparameter}
\dots

\subsection*{Eksperimen E4 — Analisis Interpretabilitas}
\dots

\subsection*{Eksperimen E5 — Generalisasi Lintas Dataset}
\dots

\section{Metrik Evaluasi}
\label{subsec:bab4_metrik}
% Persamaan RMSE, MAE, dan PHM2012 scoring function (asimetris)
\dots

\section{Lingkungan Komputasi}
\label{subsec:bab4_komputasi}
% Hardware, OS, versi PyTorch / mamba-ssm / xlstm, durasi training rata-rata
\dots
